\section*{Chapitre \ref{ch:g5:division} --- Division et multiples}

\begin{solution}{exo:g5:division:1}
$851 \div 23$~: \omterm{def:g5:division:multiple}{multiples} de $23$~: $23, 46, 69, 92, 115, 138, 161,
184, 207$. Puis $85 \div 23$~: $3$~fois ($69$), \omterm{def:g3:division:remainder}{reste} $16$~; on
abaisse le $1$~: $161 \div 23 = 7$ exactement. \omterm{def:g3:division:remainder}{Quotient} $37$, \omterm{def:g3:division:remainder}{reste}
$0$ ($23 \times 37 = 851$).

$704 \div 32$~: $70 \div 32$~: $2$~fois ($64$), \omterm{def:g3:division:remainder}{reste} $6$~; on
abaisse le $4$~: $64 \div 32 = 2$. \omterm{def:g3:division:remainder}{Quotient} $22$, \omterm{def:g3:division:remainder}{reste} $0$.
\end{solution}

\begin{solution}{exo:g5:division:2}
$9 \div 2 = 4.5$~; \quad $27 \div 4 = 6.75$~; \quad
$33 \div 5 = 6.6$~; \quad $21 \div 8 = 2.625$.
\end{solution}

\begin{solution}{exo:g5:division:3}
$54 \div 4 = 13.5$~: chacun paie $13.50$.
\end{solution}

\begin{solution}{exo:g5:division:4}
$10 \div 3 = 3.333\dots$~: le chiffre $3$ se répète à jamais --- la
division ne s'arrête pas. \omterm{def:g4:numbers:round}{Arrondi} au centième~: $3.33$.
\end{solution}

\begin{solution}{exo:g5:division:5}
\omterm{def:g5:division:multiple}{Multiples} de $7$ jusqu'à $70$~: $7, 14, 21, 28, 35, 42, 49, 56, 63,
70$. \omterm{def:g5:division:multiple}{Diviseurs} de $18$~: $1, 2, 3, 6, 9, 18$. \omterm{def:g5:division:multiple}{Diviseurs} de $24$~:
$1, 2, 3, 4, 6, 8, 12, 24$.
\end{solution}

\begin{solution}{exo:g5:division:6}
$470$~: par $2$ (se termine par $0$), par $5$, par $10$~; \omterm{def:g1:addition:def}{somme} des
chiffres $11$~: pas par $3$.

$735$~: se termine par $5$~: par $5$, pas par $2$ ni $10$~; \omterm{def:g1:addition:def}{somme}
des chiffres $15$~: par $3$.

$8\,001$~: \omterm{def:g3:numbers:evenodd}{impair}~; \omterm{def:g1:addition:def}{somme} des chiffres $9$~: par $3$ seulement.

$1\,314$~: \omterm{def:g3:numbers:evenodd}{pair}~; \omterm{def:g1:addition:def}{somme} des chiffres $9$~: par $2$ et par $3$ (donc
par $6$), pas par $5$.
\end{solution}

\begin{solution}{exo:g5:division:7}
Divisible par $3$ et $5$ signifie divisible par $15$~: entre $60$
et $80$, les nombres $60$, $75$.
\end{solution}

\begin{solution}{exo:g5:division:8}
$7.2 \div 6 = 1.2$~: chaque morceau mesure $1.2$~m (vérif.~:
$1.2 \times 6 = 7.2$).
\end{solution}

\begin{solution}{exo:g5:division:9}
$1\,000 = 24 \times 41 + 16$~: quarante et un sacs pleins, $16$
billes restantes.
\end{solution}

\begin{solution}{exo:g5:division:10}
$7 \div 5 = 1.4$~: chacun des $5$ amis obtient $1.4$ barres ---
possible en découpant les barres en dixièmes. Pour $5 \div 7$~: la
division donne $0.714285\dots$, qui ne s'arrête jamais --- aucune
part décimale exacte n'existe~; seule la \omterm{def:g4:fractions:def}{fraction} $\frac57$ la
nomme exactement.
\end{solution}

\begin{solution}{exo:g5:division:11}
\omterm{def:g1:addition:def}{Somme} des chiffres~: $5 + 2 + d = 7 + d$. Divisible par $3$ quand
$7 + d$ vaut $9$, $12$ ou $15$~: $d = 2$, $5$ ou $8$. Le plus
petit est $d = 2$ (donnant $522 = 3 \times 174$).
\end{solution}
